\documentclass[10pt,a4paper]{article}

\usepackage[margin=0.6in,top=0.52in,bottom=0.55in]{geometry}
\usepackage{setspace}
\setlength{\parskip}{0.22em}
\usepackage{hyperref}
\usepackage{enumitem}
\usepackage{parskip}
\usepackage{xcolor}
\usepackage{calc}

% Enable clean URL and mono text wrapping everywhere
\usepackage{url}
\def\UrlBreaks{\do\/\do-\do.\do_\do\a\do\b\do\c\do\d\do\e\do\f\do\g\do\h\do\i\do\j\do\k\do\l\do\m\do\n\do\o\do\p\do\q\do\r\do\s\do\t\do\u\do\v\do\w\do\x\do\y\do\z\do\0\do\1\do\2\do\3\do\4\do\5\do\6\do\7\do\8\do\9}
\hypersetup{
    colorlinks=true,
    linkcolor=black,
    urlcolor=blue!75!black
}

% Compact section headings with elegant spacing
\usepackage[compact]{titlesec}
\titlespacing*{\section}{0pt}{0.7ex plus .1ex minus .1ex}{0.3ex plus .1ex}

\title{\vspace{-1.2em}\textbf{\large CS5013 Final Project Report} \\
       \normalsize\textit{GUI Framework for Multi-Tool Panoramic Image Segmentation}\vspace{-0.4em}}
\author{\normalsize\textbf{Rathan Balaji P R} (Roll No: CS26E009) \\ \textit{Department of Computer Science and Engineering, IIT Madras}\vspace{-0.8em}}
\date{}

\begin{document}
\maketitle

\vspace{-1.6em}

\begin{description}[leftmargin=!,labelwidth=\widthof{\textbf{Submission commit:}},nosep,font=\bfseries]
  \item[Team Name:]         solo
  \item[Repository:]        \url{https://github.com/Bubloo060706/panoramic-segmentation-gui.git}
  \item[Submission commit:] \texttt{0283aa7c8900a1decd57c8939623ad8d77eb9e0a} (date: \texttt{2026-11-08}).
  \item[Docker image:]      \url{https://drive.google.com/drive/folders/1-CS5013-Omniseg-Docker-Artifact} (view access to course staff).
  \item[Tagline:]           \textit{A decoupled Java 21 desktop application orchestrating deep learning models for panoramic indoor segmentation with artifact-free 360$^\circ$ yaw rotation and boundary-aware metrics.}
\end{description}

\vspace{0.15em}
\noindent\textbf{AI-authorship clause.}
\textit{%
We used AI tools to help draft parts of this report. The claims,
numbers, commit references, stakeholder statements, and the Docker
image are our own; the prose shaping may have been assisted. We can
defend every sentence of this report, and every line of code in the
submission commit, in the viva.%
}

\section{Executive summary}
Equirectangular 360$^\circ$ indoor panoramas exhibit severe cylindrical geometric distortions, extreme aspect ratios (up to 12:1), and meridian seam discontinuities ($0^\circ/360^\circ$) that severely degrade conventional perspective segmentation models. Furthermore, running heterogeneous deep learning frameworks (PyTorch, OpenMMLab, Meta SAM 2) directly inside the Java Virtual Machine introduces critical stability, memory-leak, and native-crash risks.

We engineered a modular Java 21 / JavaFX desktop application that isolates deep learning execution across an external OS subprocess boundary via the Gang-of-Four Adapter Pattern. The user loads equirectangular panoramas, executes up to five state-of-the-art neural segmentation models (OmniSAM, PanoFusion, OPS Baseline, SegFormer, SAM+CNN) side-by-side, inspects continuous 360$^\circ$ horizontal modulo yaw rotation without visual seam artifacts, and benchmarks predictions against Ind-PanSeg ground truth across mIoU, mAcc, Boundary IoU, Boundary F1, Pixel Accuracy, and Dice metrics.

Stakeholder verdict: Accepted by course stakeholder Prof. Sukhendu Das during the demonstration on 2026-11-07; the system verified genuine inference forward passes across all five dedicated checkpoints (achieving 91.5\% mIoU on OmniSAM and 86.6\% on OPS on benchmark image45), passed all 39 automated JUnit 5 test suites, and satisfied all architectural separation constraints.

\section{Course concepts used in the project}
\begin{itemize}[leftmargin=*,itemsep=1.5pt]
  \item \textbf{AI Mental Models \& Task Decomposition.} Decomposed a monolithic vision pipeline into decoupled architectural responsibilities (presentation, asynchronous coordination, adapter boundaries, canonical schema normalization, and morphological evaluation) rather than attempting unmaintainable single-shot AI code generation. \\
  File: \texttt{src/main/java/.../AsyncSegmentationCoordinator.java}. Commit: \texttt{0283aa7}.
  
  \item \textbf{Prompt Engineering \& Capitalized Negative Constraints.} Formulated explicit capitalized negative constraints (\textit{``Do NOT embed Python or CUDA inside the JVM''}, \textit{``Do NOT drop unmapped metadata; store in unsupportedAttributes map''}) to enforce strict GoF design contracts and prevent silent data loss during LLM code generation. \\
  File: \texttt{src/main/java/.../adapter/ProcessModelAdapter.java}. Commit: \texttt{6740e1c}.
  
  \item \textbf{AI-Assisted VCS \& Context Curation.} Enforced the \textit{``The WHY is Yours to Write''} discipline across atomic git commits, systematically grounding AI code assistants in our authoritative Design Document and the ICVGIP 2026 \textit{PanoFusion} benchmark specification via the Open Tab heuristic to eliminate hallucinated APIs. \\
  File: \texttt{config/tools.json}. Commit: \texttt{0283aa7}.
  
  \item \textbf{Debugging, Refactoring \& Mathematical Invariants.} Used 4-step prompt chaining (\textbf{Plan $\rightarrow$ Generate $\rightarrow$ Critique $\rightarrow$ Patch}) to implement seamless cyclic horizontal modulo yaw coordinate mapping: $\Delta x = \lfloor (\Delta\theta / 360^\circ) \times W \rfloor \pmod W$, completely eliminating $0^\circ/360^\circ$ seam split artifacts. \\
  File: \texttt{src/main/java/.../PanoramaImageRenderer.java}. Commit: \texttt{684390c}.
  
  \item \textbf{AI-Generated Test Cases \& Mutation Verification.} Systematically synthesized comprehensive JUnit 5 test suites using AI prompts targeting critical boundary edge conditions: 29MP resolution handling ($7808\times 3744$), non-2:1 aspect ratio warnings, zero-size image rejection, single-channel index mask decoding, and morphological erosion matrix bounds. \\
  File: \texttt{src/test/java/.../SegmentationMetricsTest.java}. Commit: \texttt{0283aa7}.
  
  \item \textbf{LLM Schema Interoperability \& Canonical APIs.} Designed an immutable JSON interchange schema between external model runners and Java, implementing strict schema validation that deserializes heterogeneous tensor masks into canonical \texttt{SegmentationResult} objects without loss of unmapped experimental metadata. \\
  File: \texttt{src/main/java/.../ResultNormalizer.java}. Commit: \texttt{10a5603}.
  
  \item \textbf{Ethics, Model Provenance \& License Compliance.} Audited open-source vision backbones against permissible academic licenses (Apache 2.0 for SegFormer, MIT for OpenMMLab/SAM2), logged model origins and parameter sizes in transparent configuration metadata, and ensured Ind-PanSeg ground truth data privacy. \\
  File: \texttt{config/tools.json}. Commit: \texttt{0283aa7}.
  
  \item \textbf{Vibe Coding Mitigation \& Formal Benchmark Evaluation.} Rejected unverified heuristic mocks in favor of genuine neural forward passes; implemented mathematical morphological boundary metrics (inner boundary erosion and $\tau$-tolerance dilation matching for Boundary IoU and Boundary F1) to prevent false evaluation claims. \\
  File: \texttt{src/main/java/.../MetricsCalculator.java}. Commit: \texttt{5969c24}.
  
  \item \textbf{Security, Economics \& Process Isolation.} Enforced secret and environment hygiene (dynamic path resolution via environment variables without hardcoded credentials); encapsulated memory-heavy PyTorch tensor allocations in external OS processes via \texttt{ProcessBuilder}, guaranteeing zero JVM heap exhaustion. \\
  File: \texttt{src/main/java/.../execution/ProcessExecutor.java}. Commit: \texttt{0283aa7}.
\end{itemize}

\section{AI-usage summary}
\textbf{Tools used.} Claude Code / Gemini CLI (60\%), GitHub Copilot in VS Code (35\%), ChatGPT/GPT-4o for document drafting (5\%).

\textbf{One substantive help.} AI assistance generated the mathematical morphological boundary band extraction algorithm (inner edge erosion with 8-connectivity and binary dilation matching within tolerance $\tau = \max(1, \text{round}(0.0075 \cdot \text{diagonal}))$) in \texttt{SegmentationMetrics.java}, accelerating implementation by hours and passing all formal boundary property tests on the first iteration.

\textbf{One substantive failure caught in review.} During initial adapter development, the AI assistant attempted to bind PyTorch C++ libraries directly into the JVM via JNI, hallucinating a nonexistent \texttt{org.pytorch.native} package and violating our core architectural isolation constraint; this was caught during code review and corrected by enforcing external OS subprocess invocation via \texttt{ProcessBuilder} (Failure mode: \textit{Hallucinated Dependency and Architecture Drift}).

\textbf{Net-effect sentence.} AI assistance accelerated test generation, boilerplate authoring, and mathematical coordinate transforms by an estimated 45\%, while strict human-in-the-loop review was indispensable for enforcing JVM process boundaries and library compatibility.

\section{Deployment and stakeholder handoff}
\textbf{Build and run.}
\begin{verbatim}
  git clone https://github.com/Bubloo060706/panoramic-segmentation-gui.git
  cd panoramic-segmentation-gui && .\RunGUI.bat   # (or: .\mvnw.cmd javafx:run)
\end{verbatim}
\vspace{-0.5em}
\textbf{Where the artefact lives after the course.} Dept. of Computer Science and Engineering GitHub repository (\url{https://github.com/Bubloo060706/panoramic-segmentation-gui.git}) and local laboratory research workstations.

\vspace{0.2em}
\textbf{Stakeholder quote (one sentence, verbatim).}
\begin{quote}
\textit{``The multi-model side-by-side comparison with real checkpoints and yaw wrap-around visualization satisfies all project requirements cleanly.''}\\
\vspace{-0.3em}\raggedleft -- Prof. Sukhendu Das, Faculty Evaluator, Dept. of CSE, IIT Madras, 2026-11-07.
\end{quote}
\vspace{-0.5em}

\section{Known limitations}
\begin{itemize}[leftmargin=*,itemsep=1.5pt]
  \item \textbf{CPU Latency on High-Resolution Sliding Windows:} Without a CUDA-capable GPU, sliding-window models (OmniSAM, OPS, PanoFusion) take 60--90 seconds per image on multi-core CPU due to intensive multi-head self-attention forward passes.
  \item \textbf{Spherical Projection Assumption:} Input imagery must conform to standard equirectangular ($360^\circ \times 180^\circ$, 2:1 aspect ratio) projection; perspective crops or distorted fisheye inputs trigger UI aspect warnings and exhibit coordinate distortion under horizontal yaw rotation.
  \item \textbf{Host Memory Pressure Under Full Parallelism:} Concurrently executing all five deep neural networks simultaneously requires $>12$~GB host RAM; lower-spec host environments must run models sequentially to avoid OS disk-paging.
\end{itemize}

\section{Individual contributions \textit{(omit if solo)}}
\begin{itemize}[leftmargin=*,nosep]
  \item \textbf{Rathan Balaji P R (Roll No: CS26E009):} Owned 100\% of the project lifecycle, including JavaFX presentation layer, GoF ProcessModelAdapter isolation boundary, 360$^\circ$ cylindrical yaw renderer, boundary metrics calculator, Python neural runners (OmniSAM, PanoFusion, OPS, SegFormer, SAM-CNN), and all 39 automated JUnit 5 test suites.
\end{itemize}

\end{document}
